\section{A rational objective for counting matchings}
\label{d:sec-objective}

The weighted entropy estimate reduces counting to a continuous maximum.
We now replace that maximum by one computed from exact rational data.
Assume throughout this section that \(n\ge4\) is even.  Let \(N\) be the
number of perfect matchings in the input graph, with parallel edges
counted as distinct choices.  We allow the multiplicity \(k_e\) of each
unordered pair \(e\in\binom V2\) to be supplied in binary; an absent pair
has \(k_e=0\).  For a simple graph, every multiplicity is zero or one.

We work on the complete simple graph \(K_n\), writing
\[
  P=P(K_n),\qquad p=\binom n2.
\]
The objective functions below have one coordinate for each of these
\(p\) unordered pairs.  In particular, \(P\) is nonempty even when
the input graph has no perfect matching. Small positive weights on absent
pairs will encode the input support in the objective. Completion also
provides the uniform feasible mean \(u_{ij}=1/(n-1)\), whose nontrivial
odd-cut inequalities are strict. Section~\ref{d:sec-optimization} uses
that slack to repair cut violations after correcting the degree sums.

\subsection{Completing the graph with small weights}

Set
\begin{equation}
\label{d:eq-completion-weights}
\begin{split}
 b_{\max}&=\max\bigl(\{0\}\cup
       \{\lfloor\log_2 k_e\rfloor:k_e>0\}\bigr),\\
 L&=100(n+b_{\max}+1)^3,\qquad
 \beta_e=
 \begin{cases}
   \lfloor\log_2 k_e\rfloor,&k_e>0,\\
   -L,&k_e=0.
 \end{cases}
\end{split}
\end{equation}
For a positive integer \(k_e\), its binary bit length minus one equals
\(\lfloor\log_2 k_e\rfloor\).  Thus these coefficients are computed
exactly by integer arithmetic, and their numerical magnitudes are
polynomial in \(n+b_{\max}\).

Use \(Z_\beta\) and \(q^*\) from Corollary~\ref{d:coarse-weighted} for \(K_n\)
with the weights in Equation~\eqref{d:eq-completion-weights}.  A term in
\(Z_\beta\) is called valid if its matching uses only pairs with positive
input multiplicity; write \(Z_{\mathrm{good}}\) and
\(Z_{\mathrm{bad}}\) for the sums of valid and remaining terms,
respectively.

\begin{lemma}[Weighted completion]
\label{d:lem-penalized-count}
The contribution of matchings using an absent pair satisfies
\[
  0\le Z_{\mathrm{bad}}\le2^{-20n}.
\]
Moreover,
\begin{equation}
\label{d:eq-penalized-count}
\begin{aligned}
 N=0&\quad\Longrightarrow\quad \log_2 Z_\beta\le-20n,\\
 N>0&\quad\Longrightarrow\quad
       \max\{1,2^{-n/2}N\}\le Z_\beta\le2N.
\end{aligned}
\end{equation}
\end{lemma}

\begin{proof}
For each pair with \(k_e>0\), put
\(r_e=2^{\lfloor\log_2 k_e\rfloor}\).  Then
\[
  \frac{k_e}{2}\le r_e\le k_e.
\]
For a valid matching of \(K_n\), the product of its \(k_e\)'s is the
number of corresponding perfect matchings in the input multigraph.
Replacing these \(n/2\) factors by \(r_e\) and summing therefore gives
\begin{equation}
\label{d:eq-valid-rounded-count}
  2^{-n/2}N\le Z_{\mathrm{good}}\le N.
\end{equation}
If \(N>0\), there is at least one valid matching, and its rounded
weight is a positive integer.  Hence \(Z_{\mathrm{good}}\ge1\).

Every remaining matching uses a pair of weight \(2^{-L}\), and each
of its other pair weights is at most \(2^{b_{\max}}\).  There are at most
\(2^p\) matchings, since each is a subset of the \(p\) pairs.  Thus
\[
  Z_{\mathrm{bad}}\le 2^{p+b_{\max}n/2-L}\le2^{-20n}.
\]
For the last inequality, put \(t=n+b_{\max}+1\ge5\) and note that
\[
  p+\frac{b_{\max}n}{2}+20n\le t^2+20t\le100t^3=L.
\]
When \(N=0\), every term is in \(Z_{\mathrm{bad}}\); also \(Z_\beta>0\)
because \(K_n\) has a perfect matching and every assigned weight is
positive.  This proves the first implication in
Equation~\eqref{d:eq-penalized-count}.
When \(N>0\), Equation~\eqref{d:eq-valid-rounded-count} gives the lower
bounds, while
\[
  Z_\beta\le N+2^{-20n}\le2N
\]
because \(N\) is a positive integer.
\end{proof}

\subsection{Interpolating entropy on a dyadic grid}

Let \(K\) be the least nonnegative integer such that \(2^K\ge p\).
Define \(\phi(t)=-t\log_2 t\) on \([0,1]\), with \(\phi(0)=0\), and
let \(\psi\) be its continuous linear interpolant at
\[
  0,\ 2^{-K},\ 2^{-(K-1)},\ \ldots,\ 2^{-1},\ 1.
\]
The interpolant has the explicit rational formulas
\begin{equation}
\label{d:eq-dyadic-chords}
 \psi(t)=
 \begin{cases}
   Kt,&0\le t\le2^{-K},\\
   (k-1)t+2^{-k},
     &2^{-(k+1)}\le t\le2^{-k},\quad 0\le k<K.
 \end{cases}
\end{equation}
The formulas agree at common endpoints.

\begin{lemma}[Dyadic entropy approximation]
\label{d:lem-dyadic-entropy}
The function \(\psi\) is concave, with integer slopes of absolute value
at most \(K\).  For every \(x\in P\),
\begin{equation}
\label{d:eq-dyadic-entropy-error}
  0\le h_2(x)-\sum_e\psi(x_e)
     \le\frac n2+1\le n.
\end{equation}
\end{lemma}

\begin{proof}
At a positive knot, \(\phi(2^{-k})=k2^{-k}\).  The chord between
\(2^{-(k+1)}\) and \(2^{-k}\) therefore has slope \(k-1\) and
intercept \(2^{-k}\), whereas the chord from zero has slope \(K\).
This proves Equation~\eqref{d:eq-dyadic-chords}.  Reading the intervals
from left to right, the slopes are nonincreasing, so \(\psi\) is
concave.  Since \(p\ge6\), we have \(K\ge1\), and all its slopes lie
between \(-1\) and \(K\).

The concavity of \(\phi\) puts every chord below its graph.  If
\(2^{-(k+1)}\le t\le2^{-k}\), then
\[
  \psi(t)-kt=2^{-k}-t\ge0,
  \qquad \phi(t)\le(k+1)t.
\]
Consequently \(0\le\phi(t)-\psi(t)\le t\) on this interval.
On the interval adjacent to zero, write \(t=2^{-K}r\), where
\(0\le r\le1\).  Then
\[
  \phi(t)-\psi(t)=2^{-K}(-r\log_2 r)\le2^{-K}.
\]
Here the expression at \(r=0\) is interpreted by continuity; on
\((0,1]\), differentiation gives the maximum
\(1/(e\ln2)<1\).

For \(x\in P\), summing the degree equations gives
\(\sum_e x_e=n/2\).  Apply the bound \(x_e\) on intervals above
\(2^{-K}\), and the bound \(2^{-K}\) on the interval below it.
The total error is at most
\[
  \sum_{e:x_e\ge2^{-K}}x_e
     +\bigl|\{e:x_e<2^{-K}\}\bigr|\,2^{-K}
  \le \frac n2+p2^{-K}
  \le \frac n2+1
  \le n.
\]
This proves Equation~\eqref{d:eq-dyadic-entropy-error}.
\end{proof}

Define the concave function \(f\) on the entire box \([0,1]^p\) by
\begin{equation}
\label{d:eq-surrogate-objective}
  f(x)=\sum_e\psi(x_e)+\beta\cdot x,\qquad
  f^*=\max_{x\in P}f(x).
\end{equation}
Equation~\eqref{d:eq-dyadic-entropy-error} gives
\(f(x)\le h_2(x)+\beta\cdot x\le f(x)+n\) on \(P\).
Taking maxima yields
\begin{equation}
\label{d:eq-surrogate-gap}
  f^*\le q^*\le f^*+n.
\end{equation}
On each affine piece, \(f\) has integer coordinate slopes of absolute
value at most \(K+L\), since \(|\beta_e|\le L\).  All breakpoints and
intercepts in Equation~\eqref{d:eq-dyadic-chords} have denominators
dividing \(2^K<2p\).  Thus \(f\) can be evaluated exactly at rational
points using rational arithmetic.  Section~\ref{d:sec-optimization}
uses these properties to estimate \(f^*\).
